\section*{Standard generating sets and word problems}
Recall that a group $G$ with identity element $1$ is finitely generated (f.g.\ for short) if there is a finite set $\Sigma\subseteq G$ such that every element of $G$ can be written as a product of elements of $\Sigma$; this product can be formally written as a word from $\Sigma^*$. For technical reasons we assume that $1\in\Sigma$ (which is needed for padding reasons) and that for every $a\in\Sigma$ also the inverse $a^{-1}$ belongs to $\Sigma$; such a generating set $\Sigma$ is called standard. We have a natural involution on $\Sigma^*$ defined by $(a_1\cdots a_n)^{-1}=a_n^{-1}\cdots a_1^{-1}$ for $a_i\in\Sigma$ (which is the same as forming inverses in the group). For words $u,v\in\Sigma^*$ we write $u=_Gv$ if $u$ and $v$ are equal in $G$; sometimes we just say $u=v$ in $G$. The word problem for $G$, $\mathrm{WP}(G)$ for short, is the question whether $u=_G1$ holds for a given word $u\in\Sigma^*$. In this formulation the word problem depends on the generating set $\Sigma$, but it is well-known that the complexity/decidability status of the word problem does not depend on $\Sigma$.
\section*{SENS conventions (Definition 1 and Section 4)}
For group elements $g,h\in G$ or words $g,h\in\Sigma^*$ we write $g^h$ for the conjugate $h^{-1}gh$ and $[h,g]$ for the commutator $h^{-1}g^{-1}hg$. We call $g$ a $d$-fold nested commutator, if $d=0$ or $g=[h_1,h_2]$ for $(d-1)$-fold nested commutators $h_1,h_2$. A subquotient of a group $G$ is a quotient of a subgroup of $G$.
\begin{definition}\label{def:SENS}		
  We call a group $G$ with the finite standard generating set $\Sigma$
  \emph{strongly efficiently non-solvable} \emph{(SENS)} if for every
  $d \in \N$ there is a collection of $2^{d+1}-1$ elements
  $g_{d,v} \in \Sigma^*$ for $v \in \oneset{0,1}^{\leq d}$ such that
  \begin{enumerate}[label=(\alph*)]
  \item\label{SENSa} there is some constant $\mu \in \N$ with $\abs{g_{d,v}} = 2^{\mu d}$ for all $v \in \oneset{0,1}^{d}$,
  \item\label{SENSb} $g_{d,v} =	\bigl[ g_{d,v0},\,  g_{d,v1} \bigr]$ for all $v \in \oneset{0,1}^{< d}$ (here we take the commutator of words),
  \item\label{SENSc} $g_{d,\eps} \neq 1$ in $G$.
  \end{enumerate}
   The group $G$ is called \emph{uniformly strongly efficiently
    non-solvable} if, moreover,
  \begin{enumerate}[label=(\alph*)]
    \setcounter{enumi}{3}
  \item\label{SENSu} given $v\in \oneset{0,1}^d$, a number $i$ encoded in binary with $\mu d$ bits, and $a \in \Sigma$ one can decide in $\DLINTIME$ whether the
  	$i$-th letter of $g_{d,v}$ is $a$.
  \end{enumerate} 
\end{definition}
Here are some simple observations:
\begin{itemize}
\item A strongly efficiently non-solvable group clearly cannot
be solvable, so the above terminology makes sense.
\item If one can find suitable $g_{d,v}$ of length at most $2^{\mu d}$,  then these words can always be padded to length $2^{\mu d}$ thanks to the padding letter $1$.
\item  It suffices to specify $g_{d,v}$ for $v \in \oneset{0,1}^{d}$;
  the other $g_{d,v}$ are then defined by Condition~(b).
 \item We have $\abs{g_{d,v}} = 2^{\mu d+ 2(d-\abs{v})}$ for all $v \in \oneset{0,1}^{\leq d}$. Thus, all $g_{d,v}$ have length $2^{\Oh(d)}$.
 \item Equivalently to Condition~(d), we can require that given $v\in \oneset{0,1}^d$ and a binary encoded number $i$ with $\mu d$ bits, one can compute the $i$-th letter of $g_{d,v}$ in $\DLINTIME$.
\end{itemize}
Henceforth, whenever $d$ is clear, we simply write $g_v$ instead of $g_{d,v}$.

